\section*{Introduction}

This dictionary treats Chinese characters as linguistic evidence. Its arrangement follows phonetic series, with modern Mandarin words, radical-and-stroke ordering, pinyin, and English glosses playing only secondary roles. It records the graphic and phonological relationships that the Chinese script itself preserves. Its basic unit is the \emph{xiéshēng} 諧聲 or \emph{xíngshēng} 形聲 series, a set of characters connected by a shared phonetic element. The dictionary asks what a phonetic element tells us, how much it tells us, when it tells us too little, and how the semantic and graphic additions to a phonetic base shape the later history of the series.

A reader familiar with Chinese characters already knows the surface of this problem. Characters such as 江, 河, 湖, and 海 have semantic components that suggest the domain of water, while their other components point toward pronunciation. But this familiar schoolbook formulation becomes much less simple when it is treated historically. The phonetic component of a character seldom tells us the exact pronunciation of the word that the character writes. It points instead to a range of possible pronunciations, a historically meaningful neighborhood. Sometimes the range is narrow. Sometimes it is broad. Sometimes the script preserves an earlier relationship that the modern form has obscured. Sometimes a later graph has itself become a phonetic base and must be treated as the head of a subseries. Sometimes the evidence from Middle Chinese, Old Chinese reconstruction, \emph{Shuōwén}, Schuessler's \emph{Minimal Old Chinese and Later Han Chinese} \parencite{schuessler-2009}, Karlgren's \emph{The Book of Odes} \parencite{karlgren-1950}, Karlgren's \emph{Grammata Serica Recensa} \parencite{karlgren-1957-gsr}, Baxter's \emph{A Handbook of Old Chinese Phonology} \parencite{baxter-1992}, Baxter and Sagart's \emph{Old Chinese: A New Reconstruction} \parencite{baxter-sagart-2014}, and modern databases will not all agree. A useful dictionary must display these complications.

The immediate purpose of the present work is therefore practical. It makes the analysis of Chinese phonetic series visible in a compact form. Each entry gathers characters that belong together graphically and phonologically. It gives the Schuessler series number, the head graph, an abstract transcription of the series or subseries, the character-level semantic analysis, the Middle Chinese reading, and source identifiers such as GSR numbers when available. The broader purpose is methodological. The dictionary rejects the idea that early Chinese writing should be passed through the filter of modern standard character forms before it can be studied. It also rejects the opposite temptation, namely to dissolve the graphic system into phonetic reconstruction alone. The Chinese script works as a system of graphic signs in which phonetic, semantic, and historical information is unevenly but often richly encoded.

The form of the dictionary follows from this aim. Entries are arranged in the order of Schuessler's GSC series. The first number of an entry, as in \texttt{01-01} or \texttt{35-04}, identifies the rhyme section; the second identifies the series within that section. These numbers are not GSR item numbers and are not local entry numbers invented for this book. They are a bridge to Schuessler's organization of the material and to the broader tradition of Chinese historical phonology. Within an entry, indentation and nesting matter. A nested group has analytic force. It often indicates that a member of the broad series has become a new phonetic base, or that a phonological distinction justifies recognizing a narrower subseries. The hierarchy is part of the analysis.

\subsection*{Transcription, normalization, and the problem of \emph{kǎishū}}

The standard way of presenting early Chinese texts has long been to identify each early graph with a \emph{kǎishū} 楷書 equivalent and then to read the resulting string as Classical Chinese. This practice is useful in many ordinary editorial contexts, but it is methodologically dangerous when treated as linguistic evidence. It assumes a degree of continuity between ancient script, later standard character form, and modern language that requires evidence. It also transforms a graph produced in one orthographic and linguistic environment into a graph belonging to another. The result is an interpretive normalization, even when it resembles transcription.

A transcription may serve several purposes. It may present a text typographically in another script. It may reveal the structure of the original writing system more explicitly than the native script does. It may also regularize the idiosyncrasies of a particular written token in order to represent the language system implied by a text. These purposes often conflict. The current convention of replacing early Chinese characters with \emph{kǎishū} forms serves none of them very well. It does not put the text into a genuinely accessible non-Chinese script. It does not make the structure of the early script clearer. It often collapses distinctions that were visible in the source or imposes distinctions that were not there. Most importantly, it gives the modern standard character a false authority as the natural destination of the earlier graph.

The danger can be seen by analogy. An edition of Old English that replaced Old English word forms with modern English lookalikes would produce no gain in fidelity. It would produce an unstable hybrid, with modern spelling and modern morphemes superimposed on an earlier linguistic structure. In the same way, an early Chinese manuscript or bronze inscription belongs to the language and writing system of its own time. The editor may of course gloss it, translate it, normalize it, and compare it with later character forms. But this normalization must not be mistaken for the object itself. The critique of \emph{kǎishū} normalization follows Hill \parencite{hill-2015-transcription}, with related concerns in Takashima \parencite{takashima-2000}, Richter \parencite{richter-2003}, and Xing \parencite{xing-2005}.

The point extends beyond pre-Qin manuscripts. Even the \emph{kǎishū} script itself is historically layered. Some graphic relationships that are visible in bronze or manuscript forms are hidden in later standard forms; other apparent standard-form relationships are secondary. A dictionary that treats \emph{kǎishū} graphs as self-explanatory therefore risks treating the last stage of a long development as if it were the key to all earlier stages. The present dictionary instead treats each character as a graphic-phonological object. Standard forms remain useful, but they should not occupy the methodological place that belongs to the historical evidence.

The same point applies to pronunciation. A phonetic component is not a direct pronunciation guide in the modern sense. It does not tell us that every character containing it was pronounced the same way in Middle Chinese or Old Chinese. It tells us that the members of a series were, at some relevant stage, close enough in sound for the graphic system to use one as the phonetic basis of another. The dictionary's transcription must therefore be more abstract than a word-level reconstruction and more explicit than a mere character list.

\subsection*{What the phonetic transcription represents}

The central transcription in this dictionary differs from an Old Chinese reconstruction. It is an abstraction from a phonetic series. This distinction is essential. A reconstructed Old Chinese word attempts to represent the pronunciation of a particular lexical item at a particular historical stage. A series transcription asks a different question, namely what phonological information the phonetic element makes available across the series.

In many cases the answer is deliberately broad. A phonetic element may point to a place of articulation and a rhyme while leaving voicing, aspiration, tone-producing suffixes, prefixes, and type A/type B distinctions unresolved. In other cases the series itself may support a more precise transcription. The notation does not mechanically average the readings of all members. It tries to model how the graphic series functions. If one character in a series has an anomalous late reading, that reading should not necessarily redefine the whole series. If a subgroup is graphically coherent and phonologically narrower, the dictionary may recognize it as a subseries with its own transcription. If the evidence is insufficient, uncertainty should be displayed openly.

The series built on 刀 offers a simple starting point. Characters such as 召, 超, 昭, 炤, 沼, and 詔 belong to a family in which the phonetic element points toward a dental initial and the rhyme \texttt{-aw}, although the individual words differ in voicing, aspiration, and tone. A transcription of the series must therefore avoid pretending that 刀, 召, and 超 all shared exactly the same pronunciation. But it must also avoid treating their differing Middle Chinese readings as unrelated. The abstraction is the common phonetic information that made the series possible.

The same principle explains the use of subseries. If a graph formed from a phonetic base later becomes the phonetic base for a further set of characters, the dictionary should show this. The graphic history of the series has become hierarchical. For instance, one may ask whether characters built around 照 should be analyzed directly under 刀, under 召, or under 昭. Each choice says something different. If 刀 is treated as the phonetic, the analysis keeps the broadest relationship visible. If 召 or 昭 is treated as the new phonetic base, the analysis recognizes a derived subgroup. The dictionary is designed to make such decisions inspectable in the entry.

The notation uses capital letters for the main character of a series when appropriate. Derived series and subseries may receive lower-case or otherwise marked forms. Numbered subscripts, such as \texttt{pa₂} or \texttt{ka₃}, are used when several phonetic values would otherwise be written alike but must be distinguished. These subscripts are not local conveniences. They belong to the whole dictionary. Numbering must therefore remain continuous across rhyme sections and chapter breaks. If the same abstract form recurs in several places, the numbering records the dictionary-wide distinction among those forms, not the layout of a single page.

\subsection*{Semantic components and Latin labels}

Most Chinese characters are not phonetic elements alone. They are compounds, and the semantic component is also part of the writing system. The present dictionary represents semantic elements with abbreviated Latin labels. The choice of Latin is not an antiquarian ornament. It solves several methodological problems at once.

If semantic components were written in pinyin, Mandarin would be given a false epistemological priority. The reader might be tempted to think that the semantic component represents a Mandarin word. If semantic components were written as \emph{kǎishū} characters, the transcription would return to the very problem it is trying to escape. If they were written in English, they would sit awkwardly in non-English scholarship and might also encourage the reader to treat the label as a translation of the whole character. Latin abbreviations are deliberately artificial. They are labels for graphic semantic determiners; they do not translate the words written by whole characters.

Thus 水 may be represented by \texttt{aq} from \emph{aqua}, 木 by \texttt{arb} from \emph{arbor}, 言 by \texttt{dic} from \emph{dico}, 手 by \texttt{man} from \emph{manus}, 心 by \texttt{cor}, 山 by \texttt{mon}, 虫 by \texttt{serp}, 金 by \texttt{met}, and so on. The particular Latin words are less important than the methodological distinction they enforce. A semantic label says that this graphic component is present and is being analyzed as semantic. It does not say that the word written by the whole character ``means water'' or ``means tree.'' It also does not claim that a reader of the ancient text mentally translated the component into Latin. The Latin label is a modern scholarly device for keeping the analysis explicit.

The relative position of the semantic component and the phonetic component is indicated typographically. In the original proposal, conventions such as \texttt{A.B}, \texttt{A:B}, and \texttt{AˣB} distinguish whether one component stands beside, above, or around another. The current dictionary's LaTeX has developed its own compact conventions using superscripts, dots, colons, and nesting. A form such as \texttt{aq·ka} or \texttt{bamb:qay} is not decoration. It encodes a judgement about the graphic relationship between a semantic component and a phonetic base. For this reason the typography of the dictionary is also data. Automated tooling must preserve the raw LaTeX, the Unicode character, the semantic label, the phonetic value, and the source evidence separately.

The semantic label list is therefore an integral part of the book. It is not an appendix of glosses but a key to the notation. The reader should learn to read the labels as one reads abbreviations in a dictionary of Indo-European roots or in an epigraphic corpus. They are economical, conventional, and deliberately restrained.

\subsection*{Middle Chinese in Indological transcription}

The dictionary also gives Middle Chinese readings. These are written in an Indological transcription, in preference to Baxter's ASCII notation, pinyin, or IPA. The decision follows from a broader view of Sino-Tibetan philology. Many literary Sino-Tibetan languages, including Tibetan, Burmese, Newar, Lepcha, and Limbu, are written in scripts ultimately derived from Brahmi. Even languages whose scripts are not Indic, such as Chinese and Tangut, are studied in a world where Indic scripts and Indological transcription already play a central role. A consistent Indological orientation lowers the cost of comparison across the family and makes Chinese forms easier to place beside Tibetan, Burmese, Newar, and Sanskrit evidence. The same comparative problem appears in Balk and Janhunen's romanization of written Mongolian \parencite{balk-janhunen-1999}, in Coblin's work on \emph{'Phags-pa} Chinese \parencite{coblin-2001,coblin-2007}, and in Hill's own discussion of Middle Chinese transcription \parencite{hill-2023-indological}.

The use of Indological transcription is also a way of keeping transcription distinct from reconstruction. The purpose of the Middle Chinese form in this dictionary is to encode the categories of the rhyme books and rhyme tables in a clear and usable form. It is not a claim that we know the exact phonetic realization of every segment. IPA-based transcription is tempting because it looks scientific, but it is often too specific. It can make an unresolved philological category look like a settled phonetic fact. A good transcription should contain neither more nor less information than the writing system and the philological tradition provide.

The Indological system used here keeps much of Baxter's practical architecture but changes the letter values where Indological conventions make better sense. Palatal initials are written with \texttt{c}, \texttt{ch}, \texttt{j}, \texttt{ñ}, \texttt{ś}, and \texttt{ź} in place of Baxter-style \texttt{tsy}, \texttt{tsyh}, \texttt{dzy}, and so forth. Retroflex stops are written with dotted letters such as \texttt{ṭ}, \texttt{ṭh}, \texttt{ḍ}, and \texttt{ṇ}. Retroflex sibilants are written as \texttt{tṣ}, \texttt{tṣh}, \texttt{dẓ}, \texttt{ṣ}, and \texttt{ẓ}. The velar nasal is \texttt{ṅ}. Laryngeals are handled with special care: the glottal stop receives its own marker, while \texttt{h} and \texttt{ḫ} distinguish the relevant voiceless and voiced laryngeal categories. These choices make the Middle Chinese forms more compatible with the transcription of other literary Sino-Tibetan languages while preserving the categories that matter.

Medials and finals follow the same principle. Division-III syllables are marked with \texttt{i}, except where the initial or other notation already makes the category clear. Rounded syllables use \texttt{w}. The \emph{chóngniǔ} distinction is represented with \texttt{y} in a way that makes the rank-IV category visible without overburdening the spelling. Baxter's \texttt{o} for a value closer to {[}ʌ{]} is replaced by \texttt{ə}, because students and readers are much less likely to misread it as the vowel of English \emph{go}. Codas and tone markers largely follow Baxter, with rising and departing tones represented by \texttt{X} and \texttt{H}.

The Middle Chinese line in each entry is part of the evidence. It anchors the character-level evidence in a widely used philological system and gives the reader a way to compare the abstract series transcription with the attested Middle Chinese category. When a character's Middle Chinese form disagrees across sources, the disagreement should be recorded and not silently harmonized. The dictionary is useful precisely because it makes such conflicts visible.

\subsection*{Sources and workflow}

The present dictionary combines hand-edited scholarship and structured data. The older LaTeX source remains the authority for many editorial decisions. It contains the dense dictionary-style entries, rare characters, handwritten-style phonetic notation, semantic labels, pinyin comments, GSR identifiers, Middle Chinese readings, and notes. Around this source a set of scripts now extracts, imports, compares, and renders data. The aim is not to replace judgement with automation. It is to make the judgement easier to inspect, test, and extend.

The workflow begins with the Schuessler series. For each GSC ID, the system gathers evidence from the existing LaTeX, from Schuessler's series organization, from Baxter-Sagart/GSR-order reconstructions, from the Mand2MC data, from the \emph{shengfu} spreadsheet, and from cached component evidence such as Wiktionary's Han-compound analyses when these are useful. Each source is kept separate. A character may appear in one source and not another. A Middle Chinese reading may agree in substance but differ in transcription. A character may have a plausible semantic analysis in one tradition and a different one in another. The task is not to flatten these sources into a single undifferentiated fact, but to decide what belongs in the dictionary entry and to preserve enough provenance that the decision can be checked.

For each candidate character the workflow asks several questions. What is the visible semantic component? What is the phonetic component? Is the phonetic component the series head, a derived graph, or a deeper subseries graph? Does the standard form obscure an older relationship? Is the character a same-series variant? Is it an inherited node from the hand-edited LaTeX hierarchy, or a newly imported candidate? Does the evidence support a renderable root, or is a curated supplement needed? These questions are philological, but the pipeline forces them into a form that can be tested.

Once a candidate is placed in a series or subseries, the system resolves a root-like abstract transcription. This step is Old-Chinese-first, drawing on Baxter \parencite{baxter-1992} and Baxter and Sagart \parencite{baxter-sagart-2014} where useful. It takes Old Chinese evidence where available and abstracts away from features that the phonetic series does not require, including prefixes, voicing or manner contrasts when these are not series-defining, and final \texttt{-ʔ} or \texttt{-s} where these function as tone-producing morphology and not core series information. The resolved root is then combined with semantic evidence and rendered into the dictionary's transcription. A separate numbering step assigns dictionary-wide subscripts to duplicate phonetic transcriptions. Only after this global numbering has happened should the entries be rendered into pages.

This order matters. If numbering were done separately by rhyme section, the same form might receive different subscripts depending on page layout. If rendering were allowed to decide roots locally, the visual organization of the book would contaminate the analysis. The correct architecture collects evidence, resolves roots, builds semantic evidence, numbers globally, and then renders. The generated book may have title pages, parts, chapter-like rhyme headings, running headers, and page breaks, but these are layout features. They must not reset the linguistic analysis.

The pipeline is also intentionally conservative. It stops when a rendered candidate lacks a transliteration. Such a failure is not a nuisance but a useful signal. It tells us that the dictionary is about to print a character without a fully inspectable analysis. In recent expansion work, many such failures were genuine cases where the imported packet had only Mand2MC evidence, no BS-GSR row, no \emph{shengfu} row, no usable Wiktionary component analysis, and no resolvable root. In such cases the right response is not to weaken the test or invent a reading automatically. It is to decide whether a general rule is missing or whether the series requires a curated head-root supplement.

\subsection*{Some difficult cases}

The value of this dictionary lies especially in cases where a flat character list would conceal the problem. A few examples illustrate the kinds of judgement involved.

The series around 或 and 國 shows how a broad series may need to be split when the graphic and phonological evidence justify it. If 蟈 is treated as built directly on 或, then it shares the same broad \texttt{quyk} phonetic base as 蜮. But if 國 is understood as a derived phonetic that specifies a velar initial, then 蟈 belongs under a narrower \texttt{kuyk} subseries. The difference is not cosmetic. The first analysis emphasizes continuity with the larger 或 series; the second recognizes a meaningful derived phonetic base. The dictionary should allow the reader to see which decision has been made.

The relation between 袁 and 睘 presents a similar issue. Karlgren grouped them together, but their readings divide in a way that makes a split into \texttt{quan} and \texttt{quen} attractive. Schuessler also separates them into different series. Here the dictionary's organization by GSC series helps prevent a misleading flattening. It gives the broad historical comparison but lets the printed analysis follow the more precise division.

The 門 series illustrates a different problem. The series itself may not allow the vowel to be distinguished. Some members point to a value like \texttt{mən}, others to \texttt{mun}. A transcription such as \texttt{mvn} or \texttt{m-n} represents the ambiguity without forcing a false precision. Here the transcription is valuable because it leaves undecided what the evidence leaves undecided. A reader who wants the individual Middle Chinese forms can see them; the series value remains appropriately abstract.

The 口 series shows the opposite situation. If all relevant readings share aspiration and type-A behavior, the series transcription should include those facts. In such a case the abstract series value includes more than a bare velar plus rhyme; it records that the series itself supports a more specific shape. The notation must therefore be flexible enough to abstract away from some features in one series while preserving them in another.

The 隹 series shows how difficult the common-denominator approach can become. Some readings cluster around a value like \texttt{tuj}, others around a value like \texttt{qui}, while still others do not fit neatly into either. The graphic representation does not always allow the series to be subdivided cleanly. A dictionary must then choose between a broad notation, a split supported by reconstruction but not by visible graphic hierarchy, and a note that explains the remaining irregularity. This is exactly the sort of case where an explicit transcription is better than a silent list.

The 莫 series warns against treating every attested reading as equally decisive for the series abstraction. A few characters lack the expected final \texttt{-k}, but the open-syllable readings are plausibly late or analogical developments. To write the whole series as if \texttt{-k} were optional would misrepresent the historical function of the phonetic. The dictionary should preserve the evidence for 模, 謨, and 膜, but it should not let these later developments redefine the main series as though the phonetic base had always been ambiguous in that way. The same principle applies to 冪, where the use of the series for a different vowel is better understood as the press-ganging of the closest available phonetic base than as evidence for rewriting the whole series.

The 酉 series is a good example of scholarly disagreement. One analysis treats the whole series as sharing a lateral phonetic. Schuessler divides it into values such as \texttt{yu} and \texttt{tsu}. Baxter and Sagart posit still another distribution. For such disagreements, the printed entry should choose an analysis while preserving enough information that a reader can understand the choice. In cases of this kind, the introduction and notes should make clear that the transcription system is a way to exhibit analysis, not a machine for eliminating scholarly judgement.

Earlier script forms create still another class of complication. In \emph{kǎishū}, 蚩 and 志 do not transparently contain 之 as a phonetic component. In earlier forms the relationship is clearer. For a dictionary of the standard script, one might distinguish several unrelated written series; for a historically informed dictionary, one may instead group them under the older phonetic relationship. The same issue appears in the relation between 丁 and 正. If the oracle-bone forms are taken seriously, they may belong together; by the Zhou period the relationship may already have been less obvious. A transcription can therefore be period-specific. The question is not what relationship exists eternally among modern characters, but what relationship the script of a particular period made available to its users \parencite{shaughnessy-1991}.

The 又 series raises the relation between final \texttt{-k} and final \texttt{-ʔ}. A notation such as \texttt{quyx} can record a series that points to the rising-tone domain while still allowing velar-final evidence to be considered. Here the transcription is a hypothesis about the functioning of the series, not a one-to-one rendering of a single word.

The famous Chu graph discussed by Galambos \parencite{galambos-2006}, conventionally associated with 仁, shows why \emph{kǎishū} transcription is structurally inadequate. The graph's internal relationships are lost when one simply chooses a single standard equivalent. If the original reader understood the graph as ``a word pronounced like X with a semantic element Y,'' then the transcription should model that analysis. If, instead, the reader understood several internal graphic additions as productive semantic signs, the transcription should model that more complex structure. The correct transcription depends on how one understands the functioning of the script at that time and place. A modern standard character does not answer the question by itself.

The recent expansion of the dictionary has produced a more prosaic but equally important set of difficult cases. Several newly ingested series had no automatic path to a root because the packet contained only a single Mand2MC row and lacked BS-GSR, \emph{shengfu}, or usable Wiktionary component evidence. Characters such as 絼, 𢀷, 豓, and others therefore failed to render until a curated head-root supplement was supplied. These cases are useful reminders of the limits of automation. The absence of evidence is not a reason to print an empty transliteration; it is a reason to stop and ask whether the missing root can be supplied responsibly. The pipeline's failures are therefore part of the scholarly method.

\subsection*{How to read an entry}

A typical entry begins with a Schuessler number and a head graph. The number locates the series in the GSC order; the head graph identifies the visible or conventional series base. Immediately below or beside it appears the abstract phonetic transcription for the series. Characters then follow with their semantic-phonetic transcription and Middle Chinese reading. Comments may contain pinyin, Baxter-style source forms, GSR identifiers, or editorial reminders. These comments should not be dismissed as accidental. They are a compressed apparatus.

A simple character line may show a semantic label plus phonetic base, followed by an italic Middle Chinese form. If a character belongs to a subseries, it may be grouped under a larger displayed member. If a derived graph serves as a new phonetic base, the nested layout shows that relation. The indentation is therefore read analytically. Top-level members belong directly to the broad series; nested members belong to a derived subseries; deeper nesting signals further graphic or phonological subdivision.

The reader should keep three levels apart. First, there is the Chinese character as a written graph. Second, there is the abstract series transcription, which models the phonetic information made available by the graph or subgraph. Third, there is the Middle Chinese reading of the particular word or character. These three levels may agree closely, but they need not. A character may belong graphically to one series while its attested reading reflects analogy, borrowing, textual variation, or later development. Conversely, two characters may have similar Middle Chinese readings without belonging to the same graphic series. The dictionary is useful because it keeps these levels on the same page without collapsing them.

The semantic labels should also be read as graphic labels, not lexical glosses. A character marked with \texttt{aq} contains a water-related semantic component; it does not necessarily write a word meaning ``water.'' A character marked with \texttt{cor} contains 心 or a heart-related component; it does not necessarily write a word meaning ``heart.'' The label identifies the graphic semantic determiner and helps the reader reconstruct the structure of the compound.

The Middle Chinese forms are printed in the Indological system. A reader used to Baxter's notation should find the correspondences straightforward after a little practice. \texttt{j} corresponds to Baxter's \texttt{dzy}, \texttt{c} to \texttt{tsy}, \texttt{ṭ} to \texttt{tr}, \texttt{ḫ} to the voiced laryngeal, \texttt{ṅ} to the velar nasal, and so forth. Tone markers \texttt{X} and \texttt{H} retain their familiar function. The value of this notation is especially clear when Chinese forms are compared with Tibetan, Burmese, Sanskrit, or other Indologically transcribed material. The transcription looks less like an isolated Sinological code and more like part of a broader historical-philological practice.

\subsection*{Scope and limits}

The book is organized around the writing system and around Schuessler's phonetic series. It uses Old Chinese evidence, but it does not present itself as a reconstruction of Old Chinese word families. Cognacy, word-family relations, and morphological derivation are relevant, while the primary organizing principle remains the graphic and phonological structure of the script.

Meanings also play a secondary role. Semantic labels are graphic components, and the primary question is how characters are constructed and how their phonetic bases function. Glosses may be included when useful, especially from source data, but they do not drive the organization of the book.

The work provides tools for better transcription of early texts, while its entries concern characters and series, not continuous texts.

The current data model and scripts make the work more inspectable, but the analysis remains philological. The most important decisions, including where a series splits, which anomalies are late, which semantic component is operative, whether a same-series variant should inherit a root, and when a supplement is justified, require judgement. The point of the workflow is to make those judgements explicit and reproducible, not to pretend they can be eliminated.

The notation proposed here is a scholarly instrument. Other scholars may prefer different Latin labels, different conventions for graphic position, or different analyses of individual series. What matters is that the transcription should distinguish phonetic and semantic information, respect the historical stage of the script being described, and avoid smuggling modern Mandarin or modern \emph{kǎishū} assumptions into the evidence.

\subsection*{The shape of the book}

The book is divided into two main parts. The first gives the semantic component abbreviations and the conventions necessary to read the notation. This section is a key to the dictionary, not a separate essay. The second part contains the Schuessler series themselves, arranged by rhyme section. Each rhyme section begins on a new right-hand page and carries a heading with the section number and rhyme label. Within each section, entries follow the GSC order. The layout is compact, often two-column, because the material is dense and because the relationships among characters are easier to see when series members remain close to each other on the page.

This introduction belongs before both parts. Its task is to explain why the dictionary looks the way it does. A reader should be prepared for Latin labels, Indological Middle Chinese, nested subseries, numbered abstract roots, and the refusal to normalize everything into modern standard characters. Each of these features follows from the same principle. The dictionary treats Chinese characters as historical linguistic evidence.

The finished work will also need fuller discussion of sources, abbreviations, editorial conventions, and unresolved problems. It will need examples showing how to move from a manuscript graph, a standard character, or a source row into the dictionary's analysis. It will need an explanation of the relation between the hand-edited LaTeX and the generated curation layer. But the foundation is already clear. The dictionary is an argument in the form of a reference work. It argues that the Chinese script preserves phonological structure, that this structure can be represented more explicitly than traditional \emph{kǎishū} transcription allows, and that doing so will make the study of early Chinese language and literature more precise.
